% Generated draft: requires mathematical and source review.
\documentclass[11pt]{article}
\usepackage[margin=25mm]{geometry}
\usepackage{fontspec}
\setmainfont{Latin Modern Roman}
\usepackage{amsmath,amssymb,mathtools,unicode-math}
\setmathfont{Latin Modern Math}
\usepackage{enumitem,needspace,xurl,hyperref}
\setlength{\parindent}{0pt}
\setlength{\parskip}{5pt}
\setlength{\emergencystretch}{3em}
\newcommand{\N}{\mathbb{N}}
\newcommand{\Z}{\mathbb{Z}}
\newcommand{\Q}{\mathbb{Q}}
\newcommand{\R}{\mathbb{R}}
\newcommand{\C}{\mathbb{C}}
\newcommand{\F}{\mathbb{F}}
\newcommand{\A}{\mathbb{A}}
\newcommand{\PP}{\mathbb{P}}
\newcommand{\E}{\mathbb{E}}
\newcommand{\eps}{\varepsilon}
\newcommand{\dd}{\,\mathrm{d}}
\begin{document}
\section*{UnsolvedMath ID: 30004397}
\textbf{Source problem code:} catalogue entry
\subsection{Definitions and mathematical statement}
Let \(G\) be a group, and let \(B\) and \(D\) denote, respectively, the Bestvina complex and the Davis complex associated with the same strictly developable simple complex of groups, with their associated \(G\)-actions. The supplied record does not specify the indexing data, local groups, construction choices, or further hypotheses on this complex of groups. Thus the scope of the following questions remains subject to clarification from the cited source.

A metric \(d_B\) on \(B\) is \(G\)-invariant if
\[
 d_B(gx,gy)=d_B(x,y)\qquad(g\in G,\ x,y\in B).
\]
A geodesic metric space is \(\operatorname{CAT}(0)\) if, for every geodesic triangle and every pair of points on its sides, the distance between those points is at most the distance between the corresponding points of a Euclidean comparison triangle.

\textbf{Question 1.} For every instance of the construction intended by the source, does \(B\) admit a \(G\)-invariant \(\operatorname{CAT}(0)\) metric inducing its complex topology? Compatibility with the complex topology is the interpretation of ``a metric on the complex'' used here; the record does not specify completeness or any prescribed piecewise geometric structure.

For a specified \(G\)-equivariant topological embedding \(i:B\hookrightarrow D\), say that \(i(B)\) is a \(G\)-equivariant deformation retract of \(D\) if there are continuous maps
\[
 r:D\longrightarrow B,\qquad H:D\times[0,1]\longrightarrow D
\]
such that, for all \(b\in B\), \(x\in D\), \(g\in G\), and \(t\in[0,1]\),
\[
\begin{gathered}
 r(i(b))=b,\qquad r(gx)=g\,r(x),\\
 H(x,0)=x,\qquad H(x,1)=i(r(x)),\\
 H(i(b),t)=i(b),\qquad H(gx,t)=g\,H(x,t).
\end{gathered}
\]
Here \(G\) acts trivially on \([0,1]\).

\textbf{Question 2.} For every intended instance, is the Bestvina complex a \(G\)-equivariant deformation retract of the associated Davis complex in this sense? If a distinguished inclusion \(i\) belongs to the intended construction, the question concerns that inclusion. If none is prescribed, an existential interpretation additionally asks for such an embedding \(i\). The supplied record does not distinguish these interpretations, and this qualification remains unresolved. The pointwise fixing condition above uses the strong deformation-retract convention; the record does not explicitly specify its convention.

The two questions are retained together as in the supplied extraction. No implication between them, affirmative answer, or mathematical verification is asserted.
\subsection{Short English statement}
For the associated Bestvina and Davis complexes, ask whether the Bestvina complex always supports a \(G\)-invariant \(\operatorname{CAT}(0)\) metric and whether it is a \(G\)-equivariant deformation retract of the Davis complex. The construction hypotheses, intended inclusion, and metric and deformation-retract conventions require clarification. This is a statement-only draft requiring human mathematical review.
\subsection{Sources}
\begin{itemize}
\item Supplied local UnsolvedMath registry record, permanent ID \(30004397\), source problem code \texttt{catalogue entry}, in the Oberwolfach Reports Open Problems collection. No named proposer is supplied; the recorded proposed year is 2020.
\item Supplied report citation: \emph{Manifolds and Groups} (2021), \emph{Oberwolfach Reports}, pp.\ 417--463, DOI \texttt{10.4171/owr/2020/7}. \url{https://doi.org/10.4171/owr/2020/7}. The record attributes the problem to workshop open-problem discussions documented in Oberwolfach Reports.
\item Supplied literature reference: Petrosyan--Prytula 2020. \url{https://doi.org/10.4171/GGD/581}.
\item Supplied literature reference: Bestvina-complex developments 2024. \url{https://msp.org/agt/2024/24-2/agt-v24-n2-p12-p.pdf}.
\item Literature assessment date: \texttt{2026-08-22}. This is supplied dated curation, not independently verified here. The supplied assessment labels the general metric and deformation-retract questions unsettled and says that recent work still describes the metric question as open beyond dimension two. The background expressly requests independent verification of time-sensitive status and cited claims.
\item The supplied registry labels the problem \texttt{open} and its statement \texttt{corrected\_verified}. Its verification note reports that the official report was checked and that the two questions were retained together in the canonical extraction. These are supplied registry assertions, not independent source verification or mathematical review performed here.
\end{itemize}
\end{document}
